% !TeX program = lualatex
% !TeX root = predavanje10.tex
% ============================================================
% Zajednički preamble za sva predavanja RWA kolegija
% Uključiti s: % ============================================================
% Zajednički preamble za sva predavanja RWA kolegija
% Uključiti s: % ============================================================
% Zajednički preamble za sva predavanja RWA kolegija
% Uključiti s: \input{../common/preamble.tex}
% ============================================================

\documentclass[aspectratio=149,11pt]{beamer}
\usetheme{CambridgeUS}

\usepackage{fontspec}              % radi samo sa XeLaTeX/LuaLaTeX
\setsansfont{Latin Modern Sans}    % sans-serif font
\setmonofont{Latin Modern Mono}    % monospace font za kod
\usepackage[croatian]{babel}
\usepackage{csquotes}              % \enquote{} → hrvatski navodnici „..."
\usepackage{lmodern}
\usepackage{booktabs}
\usepackage{listings}
\usepackage{xcolor}
\usepackage{hyperref}

% ---------- Autor i logo ----------
\author{mag.ing. Josip Torić (SIT UNIZD)}
\titlegraphic{\includegraphics[height=2.5cm]{../common/SIT_logo_tekst1.pdf}}

% ---------- Listings ----------
\definecolor{codebg}{RGB}{248,248,248}
\definecolor{codeframeclr}{RGB}{220,220,220}
\lstdefinelanguage{bash}{
  morekeywords={sudo,apt,brew,pip,pipx,uv,winget,python,git,flask,source,deactivate,code,gunicorn,docker,compose,heroku,render,railway,node,npm,pnpm,yarn,uvicorn,alembic,pytest,ruff,black,precommit},
  sensitive=false,
  morecomment=[l]{\#}
}
\lstdefinelanguage{python}{
  morekeywords={from,import,def,class,return,if,elif,else,for,in,while,try,except,raise,with,as,not,and,or,is,lambda,yield,True,False,None,await,async},
  sensitive=true,
  morecomment=[l]{\#},
  morestring=[b]',
  morestring=[b]"
}
\lstdefinelanguage{sql}{
  morekeywords={SELECT,FROM,WHERE,INSERT,INTO,VALUES,UPDATE,SET,DELETE,CREATE,TABLE,PRIMARY,KEY,FOREIGN,REFERENCES,UNIQUE,NOT,NULL,INDEX,ON,ALTER,DROP,CONSTRAINT},
  sensitive=false,
  morecomment=[l]{--}
}
\lstdefinelanguage{json}{
  morecomment=[l]{//},
  morestring=[b]",
  showstringspaces=false
}
\lstdefinelanguage{typescript}{
  morekeywords={import,export,from,const,let,var,function,return,if,else,for,of,in,while,switch,case,break,async,await,try,catch,finally,throw,new,class,this,extends,implements,interface,type,enum,readonly,public,private,protected,static,get,set,true,false,null,undefined,void,string,number,boolean,any,unknown,never,typeof,keyof,as},
  sensitive=true,
  morecomment=[l]{//},
  morecomment=[s]{/*}{*/},
  morestring=[b]',
  morestring=[b]",
  morestring=[b]`
}
\lstalias{javascript}{typescript}
\lstset{
  backgroundcolor=\color{codebg},
  basicstyle=\ttfamily\footnotesize,
  frame=single,
  rulecolor=\color{codeframeclr},
  frameround=tttt,
  breaklines=true,
  tabsize=2,
  showstringspaces=false,
  upquote=true,
  columns=fullflexible
}

% ---------- Helper macros ----------
\newcommand{\codefile}[1]{\texttt{#1}}
\newcommand{\kbd}[1]{\texttt{#1}}

% ---------- Code frame environment ----------
% Automatski dodaje fragile i shrink=5, koristi se umjesto \begin{frame}
% Primjer: \begin{codeframe}{Naslov} ... \end{codeframe}
% Primjer s opcijama: \begin{codeframe}[shrink=15]{Naslov} ... \end{codeframe}
\newenvironment{codeframe}[2][]
  {\begin{frame}[fragile,shrink=5,#1]{#2}}
  {\end{frame}}

% ============================================================

\documentclass[aspectratio=149,11pt]{beamer}
\usetheme{CambridgeUS}

\usepackage{fontspec}              % radi samo sa XeLaTeX/LuaLaTeX
\setsansfont{Latin Modern Sans}    % sans-serif font
\setmonofont{Latin Modern Mono}    % monospace font za kod
\usepackage[croatian]{babel}
\usepackage{csquotes}              % \enquote{} → hrvatski navodnici „..."
\usepackage{lmodern}
\usepackage{booktabs}
\usepackage{listings}
\usepackage{xcolor}
\usepackage{hyperref}

% ---------- Autor i logo ----------
\author{mag.ing. Josip Torić (SIT UNIZD)}
\titlegraphic{\includegraphics[height=2.5cm]{../common/SIT_logo_tekst1.pdf}}

% ---------- Listings ----------
\definecolor{codebg}{RGB}{248,248,248}
\definecolor{codeframeclr}{RGB}{220,220,220}
\lstdefinelanguage{bash}{
  morekeywords={sudo,apt,brew,pip,pipx,uv,winget,python,git,flask,source,deactivate,code,gunicorn,docker,compose,heroku,render,railway,node,npm,pnpm,yarn,uvicorn,alembic,pytest,ruff,black,precommit},
  sensitive=false,
  morecomment=[l]{\#}
}
\lstdefinelanguage{python}{
  morekeywords={from,import,def,class,return,if,elif,else,for,in,while,try,except,raise,with,as,not,and,or,is,lambda,yield,True,False,None,await,async},
  sensitive=true,
  morecomment=[l]{\#},
  morestring=[b]',
  morestring=[b]"
}
\lstdefinelanguage{sql}{
  morekeywords={SELECT,FROM,WHERE,INSERT,INTO,VALUES,UPDATE,SET,DELETE,CREATE,TABLE,PRIMARY,KEY,FOREIGN,REFERENCES,UNIQUE,NOT,NULL,INDEX,ON,ALTER,DROP,CONSTRAINT},
  sensitive=false,
  morecomment=[l]{--}
}
\lstdefinelanguage{json}{
  morecomment=[l]{//},
  morestring=[b]",
  showstringspaces=false
}
\lstdefinelanguage{typescript}{
  morekeywords={import,export,from,const,let,var,function,return,if,else,for,of,in,while,switch,case,break,async,await,try,catch,finally,throw,new,class,this,extends,implements,interface,type,enum,readonly,public,private,protected,static,get,set,true,false,null,undefined,void,string,number,boolean,any,unknown,never,typeof,keyof,as},
  sensitive=true,
  morecomment=[l]{//},
  morecomment=[s]{/*}{*/},
  morestring=[b]',
  morestring=[b]",
  morestring=[b]`
}
\lstalias{javascript}{typescript}
\lstset{
  backgroundcolor=\color{codebg},
  basicstyle=\ttfamily\footnotesize,
  frame=single,
  rulecolor=\color{codeframeclr},
  frameround=tttt,
  breaklines=true,
  tabsize=2,
  showstringspaces=false,
  upquote=true,
  columns=fullflexible
}

% ---------- Helper macros ----------
\newcommand{\codefile}[1]{\texttt{#1}}
\newcommand{\kbd}[1]{\texttt{#1}}

% ---------- Code frame environment ----------
% Automatski dodaje fragile i shrink=5, koristi se umjesto \begin{frame}
% Primjer: \begin{codeframe}{Naslov} ... \end{codeframe}
% Primjer s opcijama: \begin{codeframe}[shrink=15]{Naslov} ... \end{codeframe}
\newenvironment{codeframe}[2][]
  {\begin{frame}[fragile,shrink=5,#1]{#2}}
  {\end{frame}}

% ============================================================

\documentclass[aspectratio=149,11pt]{beamer}
\usetheme{CambridgeUS}

\usepackage{fontspec}              % radi samo sa XeLaTeX/LuaLaTeX
\setsansfont{Latin Modern Sans}    % sans-serif font
\setmonofont{Latin Modern Mono}    % monospace font za kod
\usepackage[croatian]{babel}
\usepackage{csquotes}              % \enquote{} → hrvatski navodnici „..."
\usepackage{lmodern}
\usepackage{booktabs}
\usepackage{listings}
\usepackage{xcolor}
\usepackage{hyperref}

% ---------- Autor i logo ----------
\author{mag.ing. Josip Torić (SIT UNIZD)}
\titlegraphic{\includegraphics[height=2.5cm]{../common/SIT_logo_tekst1.pdf}}

% ---------- Listings ----------
\definecolor{codebg}{RGB}{248,248,248}
\definecolor{codeframeclr}{RGB}{220,220,220}
\lstdefinelanguage{bash}{
  morekeywords={sudo,apt,brew,pip,pipx,uv,winget,python,git,flask,source,deactivate,code,gunicorn,docker,compose,heroku,render,railway,node,npm,pnpm,yarn,uvicorn,alembic,pytest,ruff,black,precommit},
  sensitive=false,
  morecomment=[l]{\#}
}
\lstdefinelanguage{python}{
  morekeywords={from,import,def,class,return,if,elif,else,for,in,while,try,except,raise,with,as,not,and,or,is,lambda,yield,True,False,None,await,async},
  sensitive=true,
  morecomment=[l]{\#},
  morestring=[b]',
  morestring=[b]"
}
\lstdefinelanguage{sql}{
  morekeywords={SELECT,FROM,WHERE,INSERT,INTO,VALUES,UPDATE,SET,DELETE,CREATE,TABLE,PRIMARY,KEY,FOREIGN,REFERENCES,UNIQUE,NOT,NULL,INDEX,ON,ALTER,DROP,CONSTRAINT},
  sensitive=false,
  morecomment=[l]{--}
}
\lstdefinelanguage{json}{
  morecomment=[l]{//},
  morestring=[b]",
  showstringspaces=false
}
\lstdefinelanguage{typescript}{
  morekeywords={import,export,from,const,let,var,function,return,if,else,for,of,in,while,switch,case,break,async,await,try,catch,finally,throw,new,class,this,extends,implements,interface,type,enum,readonly,public,private,protected,static,get,set,true,false,null,undefined,void,string,number,boolean,any,unknown,never,typeof,keyof,as},
  sensitive=true,
  morecomment=[l]{//},
  morecomment=[s]{/*}{*/},
  morestring=[b]',
  morestring=[b]",
  morestring=[b]`
}
\lstalias{javascript}{typescript}
\lstset{
  backgroundcolor=\color{codebg},
  basicstyle=\ttfamily\footnotesize,
  frame=single,
  rulecolor=\color{codeframeclr},
  frameround=tttt,
  breaklines=true,
  tabsize=2,
  showstringspaces=false,
  upquote=true,
  columns=fullflexible
}

% ---------- Helper macros ----------
\newcommand{\codefile}[1]{\texttt{#1}}
\newcommand{\kbd}[1]{\texttt{#1}}

% ---------- Code frame environment ----------
% Automatski dodaje fragile i shrink=5, koristi se umjesto \begin{frame}
% Primjer: \begin{codeframe}{Naslov} ... \end{codeframe}
% Primjer s opcijama: \begin{codeframe}[shrink=15]{Naslov} ... \end{codeframe}
\newenvironment{codeframe}[2][]
  {\begin{frame}[fragile,shrink=5,#1]{#2}}
  {\end{frame}}


\title[Programiranje za web]{Programiranje za web - Predavanje 10 \\ Admin tablice, filteri i permission-aware UI}
\date{14.12.2026.}

\begin{document}

\frame{\titlepage}

% ------------------------------------------------------------
\begin{frame}{Plan predavanja}
\begin{enumerate}
  \item Admin CRUD ekrani za natjecanja
  \item Tablice s filterima i paginacijom
  \item Modali za destruktivne akcije
  \item CSV export
  \item Permission-aware UI: faze i role u frontendu
  \item Commitovi za ovo predavanje
\end{enumerate}

\vspace{0.6em}
Na kraju predavanja: admin može upravljati natjecanjima i pregledavati sve registracije s filterima i exportom.
\end{frame}

% ------------------------------------------------------------
\begin{frame}{Ishodi učenja}
Nakon ovog predavanja student može:
\begin{itemize}
  \item implementirati admin CRUD ekrane s tablicom i formom
  \item dodati filtriranje i paginaciju na tablicu
  \item koristiti modalni dijalog za potvrdu destruktivnih akcija
  \item generirati CSV export na frontendu iz podataka s API-ja
  \item prilagoditi UI ovisno o korisnikovoj roli i fazi natjecanja
\end{itemize}
\end{frame}

% ------------------------------------------------------------
\begin{frame}{Admin panel: odgovornosti}
Admin korisnik ima širi pristup od običnog korisnika. U našem sustavu:

\vspace{0.4em}
\begin{tabular}{@{}lp{8cm}@{}}
  \toprule
  \textbf{Akcija} & \textbf{Samo admin} \\
  \midrule
  Kreiranje natjecanja & Da - klubovi ne mogu kreirati natjecanja \\
  Pregled svih registracija & Da - klubovi vide samo svoje \\
  Reset lozinke kluba & Da - klubovi si ne mogu sami resetirati \\
  Kreiranje novih klubova & Da - samo admin onboarda klubove \\
  CSV export & Da - admin treba podatke za organizaciju \\
  \bottomrule
\end{tabular}

\vspace{0.6em}
Admin viewovi i club viewovi dijele isti API, ali imaju različite UI ekrane. Backend kontrolira pristup, frontend kontrolira prikaz.
\end{frame}

% ------------------------------------------------------------
\begin{frame}{Princip najmanjeg privilegija na frontendu}
Ne prikazuj korisniku akcije koje ne smije izvršiti:

\vspace{0.4em}
\begin{itemize}
  \item Club korisnik ne vidi \enquote{Novo natjecanje} gumb jer ne može kreirati natjecanja
  \item Club ne vidi admin navigaciju, ali admin vidi sve
  \item Gumbi za akcije koje ovise o fazi se skrivaju ili onemogućuju
\end{itemize}

\vspace{0.6em}
Ali: nikad se ne oslanjajte samo na frontend za sigurnost. Korisnik može:
\begin{itemize}
  \item Otvoriti DevTools i promijeniti DOM
  \item Direktno slati HTTP requestove (curl, Postman)
  \item Modificirati JavaScript u browseru
\end{itemize}

\vspace{0.4em}
Frontend sakrivanje je UX, backend provjera je sigurnost. Oboje su potrebni.
\end{frame}

% ------------------------------------------------------------
\begin{frame}{Admin vs Club viewovi}
\begin{columns}[T]
\column{0.48\textwidth}
\textbf{Admin vidi}
\begin{itemize}
  \item Sve klubove i njihove natjecatelje
  \item CRUD za natjecanja
  \item Sve registracije na natjecanje
  \item Mogućnost exporta podataka
\end{itemize}

\column{0.48\textwidth}
\textbf{Club vidi}
\begin{itemize}
  \item Samo svoj klub i svoje natjecatelje
  \item Listu natjecanja (read-only)
  \item Samo svoje registracije
  \item Akcije ovisno o fazi
\end{itemize}
\end{columns}

\vspace{0.6em}
Isti API endpointi, različiti viewovi. Backend kontrolira što korisnik \textbf{smije} - frontend kontrolira što korisnik \textbf{vidi}.
\end{frame}

% ------------------------------------------------------------
\begin{frame}{Tablica kao UI pattern}
Tablica je najčešći način prikaza strukturiranih podataka u admin panelu. Dobra tablica ima:

\vspace{0.4em}
\begin{tabular}{@{}lp{8cm}@{}}
  \toprule
  \textbf{Element} & \textbf{Svrha} \\
  \midrule
  Zaglavlje (header) & Nazivi stupaca, sortiranje klikom na header \\
  Redovi & Jedan red = jedan entitet (natjecanje, registracija) \\
  Akcije & Gumbi/linkovi u zadnjem stupcu (Uredi, Obriši, Detalji) \\
  Filteri & Search bar i dropdown iznad tablice \\
  Paginacija & Navigacija kroz stranice ispod tablice \\
  Empty state & Poruka kad nema podataka koji zadovoljavaju filtere \\
  \bottomrule
\end{tabular}
\end{frame}

% ------------------------------------------------------------
\begin{frame}{Sortiranje na frontendu}
Za male datasetove, sortiranje radimo u JavaScriptu:

\vspace{0.4em}
\begin{itemize}
  \item \texttt{Array.sort()} sortira in-place - koristite \texttt{[...array].sort()} za novi array
  \item Sortiranje stringova: \texttt{a.name.localeCompare(b.name)} za ispravno sortiranje dijakritika
  \item Sortiranje datuma: usporedba timestamp vrijednosti
  \item Multi-column sort: primarni i sekundarni kriterij
\end{itemize}

\vspace{0.6em}
UX konvencija: klik na header stupca sortira po tom stupcu. Strelica pokazuje smjer (asc/desc). Drugi klik invertira smjer. Koristite \texttt{computed} - ne mutirajte originalni array.
\end{frame}

% ------------------------------------------------------------
\begin{frame}{Filtriranje: dizajn pattern}
Filteri se kombiniraju - korisnik može istovremeno pretraživati i filtrirati po statusu:

\vspace{0.4em}
\begin{enumerate}
  \item Svaki filter je \texttt{ref()} - \texttt{search}, \texttt{statusFilter}, \texttt{categoryFilter}
  \item \texttt{computed} chain primjenjuje sve filtere redom na originalni array
  \item Paginacija se primjenjuje na filtrirani rezultat (ne na originalni)
  \item Promjena filtera resetira stranicu na 1
\end{enumerate}

\vspace{0.6em}
\begin{center}
\texttt{svi podaci $\rightarrow$ search filter $\rightarrow$ status filter $\rightarrow$ sort $\rightarrow$ paginate $\rightarrow$ prikaz}
\end{center}

\vspace{0.4em}
Svaki korak je \texttt{computed} koji ovisi o prethodnom. Vue automatski recalculira kad se bilo koji input promijeni.
\end{frame}

% ------------------------------------------------------------
\begin{frame}[fragile,shrink=5]{Admin tablica natjecanja}
\begin{lstlisting}[language=html]
<script setup lang="ts">
import { ref, onMounted } from 'vue'
import api from '@/api/client'

const competitions = ref([])
const loading = ref(true)

onMounted(async () => {
  try {
    const { data } = await api.get('/competitions')
    competitions.value = data
  } finally {
    loading.value = false
  }
})
</script>

<template>
  <h1>Natjecanja</h1>
  <RouterLink to="/admin/competitions/new" class="btn">
    Novo natjecanje
  </RouterLink>
  <div v-if="loading">Ucitavanje...</div>
  <table v-else>
    <thead>
      <tr>
        <th>Naziv</th><th>Datum</th>
        <th>Lokacija</th><th>Faza</th><th>Akcije</th>
      </tr>
    </thead>
    <tbody>
      <tr v-for="c in competitions" :key="c.id">
        <td>{{ c.name }}</td>
        <td>{{ c.date }}</td>
        <td>{{ c.location }}</td>
        <td><StatusBadge :status="c.phase" /></td>
        <td>
          <RouterLink :to="`/admin/competitions/${c.id}`">
            Detalji
          </RouterLink>
        </td>
      </tr>
    </tbody>
  </table>
</template>
\end{lstlisting}
\end{frame}

% ------------------------------------------------------------
\begin{frame}[fragile,shrink=5]{Filtriranje: search i status filter}
\begin{lstlisting}[language=html]
<script setup lang="ts">
import { ref, computed } from 'vue'

const search = ref<string>('')
const statusFilter = ref<string>('all')

const filtered = computed(() => {
  return registrations.value.filter((r) => {
    const matchesSearch =
      r.lifter_name.toLowerCase()
        .includes(search.value.toLowerCase())
    const matchesStatus =
      statusFilter.value === 'all'
      || r.status === statusFilter.value
    return matchesSearch && matchesStatus
  })
})
</script>

<template>
  <input v-model="search" placeholder="Pretrazi..." />
  <select v-model="statusFilter">
    <option value="all">Svi</option>
    <option value="active">Aktivni</option>
    <option value="withdrawn">Odjavljeni</option>
  </select>
  <!-- tablica koristi filtered umjesto registrations -->
</template>
\end{lstlisting}

\vspace{0.4em}
\texttt{computed} se automatski recalculira kad se \texttt{search} ili \texttt{statusFilter} promijeni.
\end{frame}

% ------------------------------------------------------------
\begin{frame}[fragile,shrink=5]{Paginacija na frontendu}
\begin{lstlisting}[language=html]
<script setup lang="ts">
import { ref, computed } from 'vue'

const page = ref(1)
const perPage = 10

const totalPages = computed(() =>
  Math.ceil(filtered.value.length / perPage)
)

const paginated = computed(() => {
  const start = (page.value - 1) * perPage
  return filtered.value.slice(start, start + perPage)
})
</script>

<template>
  <!-- tablica koristi paginated -->
  <div class="pagination">
    <button @click="page--" :disabled="page <= 1">
      Prethodna
    </button>
    <span>{{ page }} / {{ totalPages }}</span>
    <button @click="page++" :disabled="page >= totalPages">
      Sljedeca
    </button>
  </div>
</template>
\end{lstlisting}
\end{frame}

% ------------------------------------------------------------
\begin{frame}{Paginacija: frontend vs backend}
\begin{tabular}{@{}lp{5cm}p{5cm}@{}}
  \toprule
  & \textbf{Frontend paginacija} & \textbf{Backend paginacija} \\
  \midrule
  Kako radi & Dohvati sve, sliceaj u JS-u & Šalji \texttt{?page=2\&per\_page=10} \\
  Kad koristiti & Mali dataset (<500 redova) & Veliki dataset \\
  Prednost & Instant navigacija & Ne opterećuje memoriju \\
  Mana & Spor inicijalni load & Svaka stranica = novi request \\
  \bottomrule
\end{tabular}

\vspace{0.6em}
Za projekt s 20-50 registracija, frontend paginacija je sasvim dovoljna. Backend paginacija je implementirana na \codefile{/lifters} endpointu - frontend samo šalje query parametre.
\end{frame}

% ------------------------------------------------------------
\begin{frame}{Infinite scroll vs paginacija}
\begin{columns}[T]
\column{0.48\textwidth}
\textbf{Paginacija}
\begin{itemize}
  \item Korisnik eksplicitno traži sljedeću stranicu
  \item Zna koliko podataka ima (stranica 3/10)
  \item Može preskočiti na stranicu 7
  \item Standard za admin tablice
\end{itemize}

\column{0.48\textwidth}
\textbf{Infinite scroll}
\begin{itemize}
  \item Nove stavke se učitavaju automatski na scroll
  \item Korisnik ne zna koliko ima podataka
  \item Ne može preskočiti na sredinu
  \item Standard za feedove (Instagram, Twitter)
\end{itemize}
\end{columns}

\vspace{0.6em}
Za admin tablice koristimo paginaciju - admin treba pregled nad podacima i mogućnost skakanja na specifičnu stranicu. Infinite scroll je za beskonačne feedove poput društvenih mreža.
\end{frame}

% ------------------------------------------------------------
\begin{frame}{Responsivne tablice}
HTML tablice na mobilnim ekranima postaju nečitljive kad imaju više od 3-4 stupca:

\vspace{0.4em}
\begin{tabular}{@{}lp{8cm}@{}}
  \toprule
  \textbf{Pristup} & \textbf{Opis} \\
  \midrule
  Horizontalni scroll & \texttt{overflow-x: auto} na wrapper divu. Korisnik scrolla lijevo-desno. \\
  Sakrivanje stupaca & Manje bitni stupci se skrivaju na manjim ekranima s \texttt{@media}. \\
  Card layout & Ispod breakpointa, svaki red postaje kartica s labelom i vrijednošću. \\
  \bottomrule
\end{tabular}

\vspace{0.6em}
Za studentski projekt, horizontalni scroll je dovoljno dobar i zahtijeva jednu liniju CSS-a. Card layout je produkcijski standard, ali zahtijeva više koda.
\end{frame}

% ------------------------------------------------------------
\begin{frame}[fragile,shrink=5]{Modalni dijalog za potvrdu}
\begin{lstlisting}[language=html]
<!-- src/components/ConfirmModal.vue -->
<script setup lang="ts">
defineProps<{
  show: boolean
  title: string
  message: string
}>()
defineEmits<{
  confirm: []
  cancel: []
}>()
</script>

<template>
  <div v-if="show" class="modal-overlay" @click="$emit('cancel')">
    <div class="modal" @click.stop>
      <h3>{{ title }}</h3>
      <p>{{ message }}</p>
      <div class="modal-actions">
        <button @click="$emit('cancel')">Odustani</button>
        <button @click="$emit('confirm')" class="btn-danger">
          Potvrdi
        </button>
      </div>
    </div>
  </div>
</template>
\end{lstlisting}

\vspace{0.4em}
\texttt{@click.stop} sprečava propagaciju klika - klik na modal ne zatvara overlay.
\end{frame}

% ------------------------------------------------------------
\begin{frame}{Modali: UX smjernice}
\begin{itemize}
  \item \textbf{Kad koristiti modal}: destruktivne akcije (brisanje, odjava), potvrda akcije s posljedicama
  \item \textbf{Kad ne koristiti}: informativni sadržaj, forme za unos podataka (bolje zasebna stranica)
  \item \textbf{ESC zatvara modal} - korisnici to očekuju
  \item \textbf{Klik na overlay zatvara modal} - ali ne i klik unutar modala (\texttt{@click.stop})
  \item \textbf{Focus trap} - Tab tipka ne smije izaći iz modala dok je otvoren
\end{itemize}

\vspace{0.6em}
Modal prekida korisnikov tok rada - koristite ga samo kad je prekid opravdan (sprečavanje greške).
\end{frame}

% ------------------------------------------------------------
\begin{frame}{Destruktivne akcije: UX pattern}
Odjava natjecatelja je nepovratna akcija. UX pattern za destruktivne akcije:

\vspace{0.4em}
\begin{enumerate}
  \item Gumb za akciju je \textbf{vizualno drugačiji} - crveni gumb umjesto plavog
  \item Klik otvara \textbf{modal s pitanjem} - ne izvršava odmah
  \item Modal jasno opisuje \textbf{posljedice} - \enquote{Odjava se ne može poništiti}
  \item \textbf{Potvrdi} gumb je crveni, \textbf{Odustani} je neutralan
  \item Opcionalno: zahtijevaj unos teksta za potvrdu (npr. GitHub brisanje repozitorija)
\end{enumerate}

\vspace{0.6em}
Alternativa modalu je \textbf{undo pattern}: akcija se izvrši odmah, ali korisnik ima 5 sekundi da klikne \enquote{Poništi}. Gmail koristi ovo za brisanje mailova.
\end{frame}

% ------------------------------------------------------------
\begin{frame}{Teleport u Vueu}
Problem: modal je definiran unutar child komponente, ali treba biti renderiran na vrhu DOM stabla (izvan layouta, iznad svih elemenata).

\vspace{0.4em}
\begin{itemize}
  \item Bez Teleporta: modal je nested u parent DOM-u. CSS \texttt{z-index} i \texttt{overflow: hidden} na parentu mogu sakriti overlay.
  \item S Teleportom: modal se renderira izvan parent DOM-a, direktno u \texttt{<body>}.
\end{itemize}

\vspace{0.6em}
\texttt{<Teleport to="body">} premješta sadržaj u \texttt{<body>} element. Komponenta ostaje u Vue stablu (props, events rade normalno), ali DOM node je na drugom mjestu. Ovo rješava z-index i overflow probleme.
\end{frame}

% ------------------------------------------------------------
\begin{frame}[fragile,shrink=5]{Korištenje modala za odjavu}
\begin{lstlisting}[language=html]
<script setup lang="ts">
import { ref } from 'vue'
import ConfirmModal from '@/components/ConfirmModal.vue'

const showModal = ref(false)
const selectedReg = ref(null)

function askWithdraw(reg) {
  selectedReg.value = reg
  showModal.value = true
}

async function confirmWithdraw() {
  await api.post(
    `/competitions/${compId}/registrations/${selectedReg.value.id}/withdraw`
  )
  showModal.value = false
  await refreshData()
}
</script>

<template>
  <button @click="askWithdraw(reg)">Odjavi</button>
  <ConfirmModal :show="showModal"
    title="Odjava natjecatelja"
    message="Jeste li sigurni? Odjava se ne moze ponistiti."
    @confirm="confirmWithdraw"
    @cancel="showModal = false" />
</template>
\end{lstlisting}
\end{frame}

% ------------------------------------------------------------
\begin{frame}[fragile,shrink=5]{CSV export}
\begin{lstlisting}[language=typescript]
// src/utils/csv.ts
export function downloadCSV(
  data: Record<string, unknown>[], filename: string
) {
  if (!data.length) return

  const headers = Object.keys(data[0])
  const rows = data.map((row) =>
    headers.map((h) => `"${row[h] ?? ''}"`).join(',')
  )

  const csv = [headers.join(','), ...rows].join('\n')
  const blob = new Blob([csv], { type: 'text/csv' })
  const url = URL.createObjectURL(blob)

  const a = document.createElement('a')
  a.href = url
  a.download = filename
  a.click()
  URL.revokeObjectURL(url)
}
\end{lstlisting}

\vspace{0.4em}
Generiramo CSV na frontendu iz podataka koje već imamo. Za velike datasetove, backend bi mogao imati dedicirani export endpoint.
\end{frame}

% ------------------------------------------------------------
\begin{frame}[fragile,shrink=5]{CSV export: korištenje u viewu}
\begin{lstlisting}[language=html]
<script setup lang="ts">
import { downloadCSV } from '@/utils/csv'

function exportRegistrations() {
  const rows = registrations.value.map((r) => ({
    natjecatelj: r.lifter_name,
    kategorija: r.category,
    status: r.status,
    datum_prijave: r.registered_at,
  }))
  downloadCSV(rows, `registracije-${compId}.csv`)
}
</script>

<template>
  <button @click="exportRegistrations">
    Preuzmi CSV
  </button>
</template>
\end{lstlisting}
\end{frame}

% ------------------------------------------------------------
\begin{frame}{Permission-aware UI}
Frontend prilagođava prikaz ovisno o \textbf{roli korisnika} i \textbf{fazi natjecanja}:

\vspace{0.4em}
\begin{tabular}{@{}llll@{}}
  \toprule
  \textbf{Element} & \textbf{Admin} & \textbf{Club (OPEN)} & \textbf{Club (CLOSED)} \\
  \midrule
  Gumb \enquote{Nova prijava} & vidljiv & vidljiv & skriven \\
  Gumb \enquote{Odjavi} & vidljiv & vidljiv & skriven \\
  Gumb \enquote{Uredi kategoriju} & vidljiv & vidljiv & skriven \\
  Link \enquote{Novo natjecanje} & vidljiv & skriven & skriven \\
  CSV export & vidljiv & skriven & skriven \\
  \bottomrule
\end{tabular}

\vspace{0.6em}
Frontend sakriva gumbe za bolji UX, ali backend i dalje odbija nedozvoljene akcije. Nikad se ne oslanjajte samo na frontend za sigurnost.
\end{frame}

% ------------------------------------------------------------
\begin{frame}{Disabled vs Hidden}
Kad korisnik nema pravo na akciju, imamo dva izbora:

\vspace{0.4em}
\begin{columns}[T]
\column{0.48\textwidth}
\textbf{Hidden (v-if)}
\begin{itemize}
  \item Element ne postoji u DOM-u
  \item Korisnik ne zna da opcija postoji
  \item Čišći UI
  \item Koristiti za role-based skrivanje
\end{itemize}

\column{0.48\textwidth}
\textbf{Disabled (:disabled)}
\begin{itemize}
  \item Element je vidljiv, ali neklikabilan
  \item Korisnik zna da opcija postoji
  \item Tooltip objašnjava zašto
  \item Koristiti za phase-based ograničenja
\end{itemize}
\end{columns}

\vspace{0.6em}
Club korisnik ne vidi \enquote{Novo natjecanje} jer tu akciju nikad neće smjeti izvesti. \enquote{Nova prijava} nakon roka je disabled jer ju je ranije mogao izvesti.
\end{frame}

% ------------------------------------------------------------
\begin{frame}{Faze natjecanja na frontendu}
Backend šalje fazu natjecanja u API odgovoru. Frontend koristi tu informaciju za prilagodbu UI-ja:

\vspace{0.4em}
\begin{tabular}{@{}llll@{}}
  \toprule
  \textbf{Faza} & \textbf{Nova prijava} & \textbf{Promjena kat.} & \textbf{Odjava} \\
  \midrule
  OPEN & Omogućeno & Omogućeno & Omogućeno \\
  PRELIM\_PASSED & Disabled + tooltip & Omogućeno & Omogućeno \\
  CLOSED & Hidden & Disabled & Disabled \\
  \bottomrule
\end{tabular}

\vspace{0.6em}
Problem: korisnik može otvoriti stranicu u OPEN fazi, otići na kavu i pokušati akciju kad je već PRELIM\_PASSED. Rješenja:
\begin{itemize}
  \item Backend i dalje odbija - korisnik dobije grešku
  \item Periodički re-fetch podataka (polling)
  \item Prikazati countdown do roka na stranici
\end{itemize}
\end{frame}

% ------------------------------------------------------------
\begin{frame}{Bulk akcije u tablicama}
Admin tablica s mnogo zapisa može trebati akcije nad više redova odjednom:

\vspace{0.4em}
\begin{itemize}
  \item Checkbox u svakom redu + \enquote{Select All} u headeru
  \item Toolbar s akcijama se pojavi kad je barem jedan red selektiran
  \item Akcije: \enquote{Exportaj odabrane}, \enquote{Obriši odabrane}
  \item Prikaz broja selektiranih: \enquote{5 od 23 odabrano}
\end{itemize}

\vspace{0.6em}
Implementacija: \texttt{selectedIds = ref(new Set<number>())}. Toggle dodaje/uklanja ID. \enquote{Select All} dodaje sve ID-eve s trenutne stranice (ne sve stranice). Bulk akcija šalje listu ID-eva na backend.
\end{frame}

% ------------------------------------------------------------
\begin{frame}[fragile,shrink=5]{Disabling akcija po fazi}
\begin{lstlisting}[language=html]
<script setup lang="ts">
import { computed } from 'vue'

const canRegister = computed(() =>
  competition.value.phase === 'open'
)
const canModify = computed(() =>
  ['open', 'prelim_passed'].includes(competition.value.phase)
)
</script>

<template>
  <button @click="register" :disabled="!canRegister"
    :title="canRegister ? '' : 'Rok za prijave je istekao'">
    Nova prijava
  </button>
  <button @click="withdraw(reg)" v-if="canModify">
    Odjavi
  </button>
</template>
\end{lstlisting}

\vspace{0.4em}
Koristimo \texttt{:disabled} umjesto \texttt{v-if} za gumbe koji su vidljivi ali nedostupni - korisnik vidi da opcija postoji, ali nije dozvoljena.
\end{frame}

% ------------------------------------------------------------
\begin{frame}{Role-based UI rendering}
Tri pristupa za prikazivanje elemenata ovisno o roli:

\vspace{0.4em}
\begin{tabular}{@{}lp{8cm}@{}}
  \toprule
  \textbf{Pristup} & \textbf{Primjer} \\
  \midrule
  Direktan v-if & \texttt{v-if="auth.isAdmin"} - jednostavno, ali ponovljeno \\
  Helper funkcija & \texttt{v-if="canCreate('competition')"} - centralizirana logika \\
  Zasebni viewovi & \texttt{AdminDashboard.vue} vs \texttt{ClubDashboard.vue} \\
  \bottomrule
\end{tabular}

\vspace{0.6em}
Za naš projekt: admin i club imaju zasebne viewove jer je UI značajno drugačiji. Unutar zajedničkih komponenti koristimo \texttt{v-if="auth.isAdmin"} za manje razlike.
\end{frame}

% ------------------------------------------------------------
\begin{frame}{Audit trail: zašto pratiti akcije}
U produkcijskom sustavu treba znati tko je što napravio i kad:

\vspace{0.4em}
\begin{itemize}
  \item Tko je odjavio natjecatelja? Admin ili klub?
  \item Kad je promijenjena kategorija? Prije ili nakon prelim roka?
  \item Je li admin resetirao lozinku kluba?
\end{itemize}

\vspace{0.6em}
\begin{tabular}{@{}lp{8cm}@{}}
  \toprule
  \textbf{Razina} & \textbf{Kako} \\
  \midrule
  Backend logovi & Strukturirani logovi: user\_id, akcija, timestamp, IP \\
  Baza & \texttt{created\_at}, \texttt{updated\_at}, opcionalno \texttt{updated\_by} \\
  Frontend & Console.log za debug, ne za audit \\
  \bottomrule
\end{tabular}

\vspace{0.4em}
Za ovaj projekt imamo \texttt{registered\_at} na registraciji. U produkciji bismo dodali \texttt{withdrawn\_at}, \texttt{withdrawn\_by} i slično.
\end{frame}

% ------------------------------------------------------------
\begin{frame}{Commitovi za Predavanje 10}
\begin{itemize}
  \item \texttt{feat(web): add admin competitions CRUD screens}
  \item \texttt{feat(web): add admin registrations table with filters and pagination}
  \item \texttt{feat(web): add approve/reject UI with modal and reason}
  \item \texttt{feat(web): add CSV export button and download flow}
  \item \texttt{fix(web): disable actions based on competition phase (prelim/final/locked)}
\end{itemize}

\vspace{0.4em}
\textbf{Na satu:} gradimo admin ekrane, testiramo filtriranje, gledamo kako se UI mijenja po fazama.
\end{frame}

% ------------------------------------------------------------
\begin{frame}{Što ide sljedeće (Predavanje 11)}
\begin{itemize}
  \item Production build i environment config
  \item Deploy na Railway ili Render
  \item CORS konfiguracija
  \item CI/CD pipeline
  \item Priprema za 2. obranu
\end{itemize}

\vspace{0.8em}
\begin{center}
\Large Pitanja?
\end{center}
\end{frame}

\end{document}
